\documentclass[11pt]{book}

\usepackage[english]{babel}
\usepackage{caption}
\usepackage{cite}
\usepackage{color}
\usepackage[T1]{fontenc}
\usepackage{graphicx}
\usepackage[utf8]{inputenc}
\usepackage[hidelinks]{hyperref}
\usepackage{listings}
\usepackage{listings-rust}
\usepackage{newclude}
\usepackage{pdfpages}
\usepackage{todonotes}

\setlength{\textwidth}{146.8mm} % = 210mm - 37mm - 26.2mm
\setlength{\oddsidemargin}{11.6mm} % 37mm - 1in (from hoffset)
\setlength{\evensidemargin}{0.8mm} % = 26.2mm - 1in (from hoffset)
\setlength{\topmargin}{-2.2mm} % = 0mm -1in + 23.2mm
\setlength{\textheight}{221.9mm} % = 297mm -29.5mm -31.6mm - 14mm (12 to accomodate footline with pagenumber)
\setlength{\headheight}{14pt}
\setlength{\parindent}{0pt}

\graphicspath{{../diagram}{../graphs}}

\usepackage{setspace}
\setstretch{1.1}

\newcommand{\marios}[1]{\noindent{{\bf \fbox{MK:} {\textcolor{green}{\it#1}}}}}

\makeatletter
\setlength{\@fptop}{0pt}
\makeatother

\usepackage{fancyhdr}
\renewcommand{\sectionmark}[1]{\markright{\thesection\ #1}}
\pagestyle{fancy}
\fancyhf{}
\renewcommand{\headrulewidth}{0.4pt}
\renewcommand{\footrulewidth}{0pt}
\fancyhead[OR]{\bfseries \nouppercase{\rightmark}}
\fancyhead[EL]{\bfseries \nouppercase{\leftmark}}
\fancyfoot[EL,OR]{\thepage}
\fancypagestyle{plain}{
  \fancyhf{}
  \renewcommand{\headrulewidth}{0pt}
  \renewcommand{\footrulewidth}{0pt}
  \fancyfoot[EL,OR]{\thepage}}
\fancypagestyle{addpagenumbersforpdfimports}{
  \fancyhead{}
  \renewcommand{\headrulewidth}{0pt}
  \fancyfoot{}
  \fancyfoot[RO,LE]{\thepage}
}

\usepackage{tikz}
\usepackage[explicit]{titlesec}
\newcommand*\chapterlabel{}
% \renewcommand{\thechapter}{\Roman{chapter}}
\titleformat{\chapter}[display]  % type (section,chapter,etc...) to vary,  shape (eg display-type)
{\normalfont\bfseries\Huge} % format of the chapter
{\gdef\chapterlabel{\thechapter\ }}     % the label
{0pt} % separation between label and chapter-title
{\begin{tikzpicture}[remember picture,overlay]
    \node[yshift=-8cm] at (current page.north west)
    {\begin{tikzpicture}[remember picture, overlay]
        \draw[fill=black] (0,0) rectangle(35.5mm,15mm);
        \node[anchor=north east,yshift=-7.2cm,xshift=34mm,minimum height=30mm,inner sep=0mm] at (current page.north west)
        {\parbox[top][30mm][t]{15mm}{\raggedleft $\phantom{\textrm{l}}$\color{white}\chapterlabel}};  %the black l is just to get better base-line alingement
        \node[anchor=north west,yshift=-7.2cm,xshift=37mm,text width=\textwidth,minimum height=30mm,inner sep=0mm] at (current page.north west)
        {\parbox[top][30mm][t]{\textwidth}{\color{black}#1}};
      \end{tikzpicture}
    };
  \end{tikzpicture}
  \gdef\chapterlabel{}
} % code before the title body

\titlespacing*{\chapter}{0pt}{50pt}{30pt}
\titlespacing*{\section}{0pt}{13.2pt}{*0}  % 13.2pt is line spacing for a text with 11pt font size
\titlespacing*{\subsection}{0pt}{13.2pt}{*0}
\titlespacing*{\subsubsection}{0pt}{13.2pt}{*0}

\newcounter{myparts}
\newcommand*\partlabel{}
\titleformat{\part}[display]  % type (section,chapter,etc...) to vary,  shape (eg display-type)
{\normalfont\bfseries\Huge} % format of the part
{\gdef\partlabel{\thepart\ }}     % the label
{0pt} % separation between label and part-title
{\setlength{\unitlength}{20mm}
  \addtocounter{myparts}{1}
  \begin{tikzpicture}[remember picture,overlay]
    \node[anchor=north west,xshift=-65mm,yshift=-6.9cm-\value{myparts}*20mm] at (current page.north east) % for unknown reasons: 3mm missing -> 65 instead of 62
    {\begin{tikzpicture}[remember picture, overlay]
        \draw[fill=black] (0,0) rectangle(62mm,20mm);   % -\value{myparts}\unitlength
        \node[anchor=north west,yshift=-6.1cm-\value{myparts}*20mm,xshift=-60.5mm,minimum height=30mm,inner sep=0mm] at (current page.north east)
        {\parbox[top][30mm][t]{55mm}{\raggedright \color{white}Part \partlabel $\phantom{\textrm{l}}$}};  %the phantom l is just to get better base-line alingement
        \node[anchor=north east,yshift=-6.1cm-\value{myparts}*20mm,xshift=-63.5mm,text width=\textwidth,minimum height=30mm,inner sep=0mm] at (current page.north east)
        {\parbox[top][30mm][t]{\textwidth}{\raggedleft \color{black}#1}};
      \end{tikzpicture}
    };
  \end{tikzpicture}
  \gdef\partlabel{}
} % code before the title body

\usepackage{etoolbox}
\let\bbordermatrix\bordermatrix
\patchcmd{\bbordermatrix}{8.75}{4.75}{}{}
\patchcmd{\bbordermatrix}{\left(}{\left[}{}{}
    \patchcmd{\bbordermatrix}{\right)}{\right]}{}{}
\usepackage{amsmath}
% Fix the problem with delimiter size caused by fourier and amsmath packages.
\makeatletter
\def\resetMathstrut@{%
  \setbox\z@\hbox{%
    \mathchardef\@tempa\mathcode`\(\relax
    \def\@tempb##1"##2##3{\the\textfont"##3\char"}%
    \expandafter\@tempb\meaning\@tempa \relax
  }%
  \ht\Mathstrutbox@1.2\ht\z@ \dp\Mathstrutbox@1.2\dp\z@
}
\makeatother

\lstset{
  backgroundcolor=\color{white},
  basicstyle=\footnotesize,
  breaklines=false,
  deletekeywords= {stack},
  language=Rust,
  frame=single,
  keywordstyle=\color{green},
}

\author{Ogier Bouvier}
\date{\today}
\title{A zero-copy key-value store in Rust}

\begin{document}

\pagenumbering{gobble}

\begin{titlepage}
  \newcommand{\HRule}{\rule{\linewidth}{0.5mm}} % Defines a new command for the horizontal lines, change thickness here
  \center

  \LARGE \textsc{École Polytechnique Fédérale de Lausanne}
  \vspace{1cm}

  \LARGE Master Thesis
  \vspace{2cm}

  {
    \HRule \\[0.5cm]
    \Huge A Zero Copy Key-Value Store in Rust}
  \HRule \\[1cm]

  \begin{minipage}[t]{0.4\textwidth}
    \centering
    \large
    {\large \today}
  \end{minipage}\\[1cm]

  \vspace{2cm}

  \begin{tabular}{c c}
    \begin{minipage}[t]{0.4\textwidth}
      \begin{flushleft} \large
        \textbf{Author} \\
        Ogier \textsc{Bouvier}\\
        \textit{ogier.bouvier@epfl.ch}
      \end{flushleft}
    \end{minipage}

    \begin{minipage}[t]{0.4\textwidth}
      \begin{flushright} \large
        \textbf{Supervisors} \\
        Prof. Edouard \textsc{Bugnion}\\
        DCSL | EPFL \\
        \textit{edouard.bugnion@epfl.ch} \\
        Marios \textsc{Kogias}\\
        DCSL | EPFL \\
        \textit{marios.kogias@epfl.ch}
      \end{flushright}
    \end{minipage}
  \end{tabular}

  \begin{figure}[h]
    \centering
    \includegraphics[width=0.5\textwidth]{../diagram/epfl.pdf}
  \end{figure}

\end{titlepage}

\pagenumbering{arabic}
\setcounter{page}{1}

\chapter*{Acknowledgements}

Marios Kogias, for his invaluable help during the project and
proof-reading of this paper

\chapter*{Abstract}
To provide a key-value store with consistent latency independent from
keys and values sizes we need to eliminate copies from the key-value
store's datapath while still ensuring reliable delivery of responses
to requests. We leverage several feature of the Rust programming
language to achieve zero-copy reliable delivery of key-value store
responses through the R2P2 protocol while providing a much higher
level of safety than these type of systems usually do.


% The state of the art today for high-speed networking involves kernel
% bypassing techniques making use of software such as DPDK to reduce the
% overhead typically associated with kernel based networking. Such
% software is usually much faster than traditional kernel networking,
% mostly due to avoidance of context switches and data copies between
% kernel and user space, especially when processing lots of small
% packets which often happens in the case of key-value stores such as
% memcached. Nonetheless these systems are far from perfect. First,
% having all the  networking code and the user application inside the
% same address space means that user code can corrupt the networking
% stack very easily. Secondly, development of such applications is also
% much more cumbersome than traditional kernel based approaches, indeed
% having the user program in the same address space as the networking
% code creates the need for it to be aware of low level networking
% details that traditional applications don't need to worry about.The
% necessity to handle this kind of networking details make kernel bypass
% application much more time consuming and error-prone to develop than
% regular socket based ones. Another problem these kind of systems
% encounter is their inability to ensure safety, since contrary to
% kernel based systems the networking stack and the user application
% share a virtual address space languages like C do not provide
% sufficient safety guarantees when data is shared between multiple
% threads.

%%% Local Variables:
%%% mode: latex
%%% TeX-master: "master"
%%% End:


\tableofcontents

To provide a key-value store with consistent latency independent from
keys and values sizes we need to eliminate copies from the key-value
store's datapath while still ensuring reliable delivery of responses
to requests. We leverage several feature of the Rust programming
language to achieve zero-copy reliable delivery of key-value store
responses through the R2P2 protocol.
\todo{performance numbers}

% The state of the art today for high-speed networking involves kernel
% bypassing techniques making use of software such as DPDK to reduce the
% overhead typically associated with kernel based networking. Such
% software is usually much faster than traditional kernel networking,
% mostly due to avoidance of context switches and data copies between
% kernel and user space, especially when processing lots of small
% packets which often happens in the case of key-value stores such as
% memcached. Nonetheless these systems are far from perfect. First,
% having all the  networking code and the user application inside the
% same address space means that user code can corrupt the networking
% stack very easily. Secondly, development of such applications is also
% much more cumbersome than traditional kernel based approaches, indeed
% having the user program in the same address space as the networking
% code creates the need for it to be aware of low level networking
% details that traditional applications don't need to worry about.The
% necessity to handle this kind of networking details make kernel bypass
% application much more time consuming and error-prone to develop than
% regular socket based ones. Another problem these kind of systems
% encounter is their inability to ensure safety, since contrary to
% kernel based systems the networking stack and the user application
% share a virtual address space languages like C do not provide
% sufficient safety guarantees when data is shared between multiple
% threads.

\marios{You need a flat abstract without subsection where you clearly
  state the problem. "We need to eliminate copies in a KV-store's
  datapath. Thus we use language features to reason about data
  onwership." You mention that we use Rust for that without any
  details. You mention that the solution is for reliable transport
  protocols and we use R2P2 without getting into details. And then you
  state some basic performance numbers.}

%%% Local Variables:
%%% mode: latex
%%% TeX-master: "master"
%%% End:

\marios{In this section we need to show you read and understood the
  area. This is where most of your references will come from.}
\section{Kernel-bypassing}

Today's high end NICs support speeds that make traditional kernel
networking too slow to fully exploit their capacity. The overhead of
context switches and copying is just too much when treating millions
of packets par second. Moreover the Linux kernel TCP stack is too
general for specific applications, having to handle so many quirks of
TCP makes it slower than could be achieve with a specifically designed
networking stack. To address these problems, a technique called kernel
bypassing was introduced allowing user space programs to speak
directly to the NIC through the PCIe bus thereby eliminating the
kernel overhead. These systems have networking performance orders of
magnitude faster than kernel based ones. DPDK is perhaps the most well
know of these kind of framework, it is developed by Intel and supports
most high end Intel NICs.

These kernel bypass systems have allowed software defined networks to
rival performance of hardware based network middleboxes through pure
software implementation. But they can also be useful for end host
networking purposes. Key-value store are a big part of most of today's
web services (Facebook, Linkedin and Amazon notably makes heavy use of
key-value stores). Applying kernel bypass techniques to key-value
store could bring huge performance improvements, by allowing a
codesign of the networking stack and the user application for our
specific purpose. MICA~\cite{mica} is an example of this, an holistic
approach to key-value stores.

However these systems also come with drawbacks, they offer no
guarantee of safety since everything runs in the same userspace
process it is easy to corrupt the networking stack from the user
application. NetBricks~\cite{netbricks} is a Rust framework built on
top of DPDK, aiming to provide isolated software defined network
functions while eliminating the need for virtualization. It builds
upon DPDK to provide a convenient Rust interface and provide safety
guarantees close to that of virtualization with a much lower
performance penalty.

\section{Transport protocols for RPCs}

The choice of transport protocol is critical to performance when
designing an efficient RPC protocol. Providing reliability is
necessary when ensuring that applications behave as expected. However
ensuring reliability usually involves creating copies of the data sent
through the network to be able to retransmit it if needed while not
sharing data with the user program. Traditional reliable networking
protocol (such as the Linux kernel's TCP stack) make heavy use of
copies to avoid this problem. However this makes the latency dependent
on the size of data therefore degrading performance, especially when
the payloads grow in size. The question is then can we provide a
reliable RPC protocol while avoiding copies altogether?

Using TCP would make such a networking stack a little bit too complex,
the constraint of avoiding data copies make delivering the TCP byte
stream a complex task. UDP on the other hand only provides a message
oriented delivery therefore being much simpler to implement. But it
does not provide any kind of reliable delivery of packets, so there is
need for need some other transport protocol on top of UDP to provide
reliable delivery of packets. R2P2, the Request Response Pair
Protocol, is a highly scalable RPC protocol, that provides
load-balanced delivery of multi-packet requests and responses on top
of UDP.

\section{Rust}

Rust is a systems programming language from Mozilla Research, stemming
from Mozilla's dissatisfaction with C++. As of today several parts of
Firefox's rendering engine, Servo, have been rewritten in Rust after
the language was deemed mature enough to do so. Rust is engineered to
provide a lot of safety guarantees, notably memory and concurrency
safety through its type system thereby eliminating a lot of common
problems in systems programming.

The strong Rust guarantees also speed up development significantly by
eliminating whole categories of bugs. Most importantly, Rust's type system
eliminates use-after-free, double-free, invalid pointers, memory leaks
and data-races. Overall, the time writing code is increasing since the
compiler will often refuse straight out to compile code that
potentially violates its guarantees. The counter-part to this is that
the time spent debugging is very much reduced. This is particularly
useful when tracking down data race bugs which are notoriously hard to
diagnose. Those kind of bugs are usually found much later in the
development cycle and are usually quite hard to reproduce since they
depend on a certain scheduling of the execution.

We make the claim that Rust~\cite{rustbook} is a better suited
language than C for kernel bypass networked applications. The safety
guarantees Rust brings are very well suited for this kind of
applications, while the performance loss (if any) is negligible. As an
additional benefit Rust allows code to be much more concise and
abstract while drastically cutting down on the memory and concurrency
bug tracking time. Its memory management model will also prove
invaluable when ensuring integrity of data in between the networking
stack and the user application.

As a proof of concept we implement a key-value store on top of
DPDK in order to prove that it is doable as well as reasonable,
performance-wise, to use Rust in this kind of scenario.


\section{Memory safety}

One of the most important feature of Rust is its guarantee of memory
safety, providing both the speed of manual memory management and the
safety of garbage collected languages.

To provide this feature, Rust introduces a new concept called
lifetimes and makes raw pointers second class citizens. These raw
pointers (such as that of C) can't be dereferenced since the compiler
can't guarantee that they are valid. Pointers are hard to validate and
verify statically mostly because of pointer arithmetic. References, on
the other hand, are guaranteed to always point to a valid memory
location. Like the references in C++, references in Rust do not allow
pointer arithmetic, they can only be created from a variable, which
ensures that the referenced memory is valid.

Lifetimes are the central part of Rust's memory safety
guarantees. They work in conjunction with references to ensure
that references are used correctly (i.e.\ not after the value they
refer to has been destroyed), since they allow the compiler to tell
when a reference might outlive the value it refers to. The combination
of lifetimes and references thus eliminate the dangling pointers
problem entirely. Lifetimes are Rust's way of eliminating the need for
garbage collection while not requiring users to manually manage
memory. A lifetime gets attached to each variable created, the
compiler then analyzes the control flow of the program to ensure that
references or values derived from that first value (such as reference
to a struct member) are not used after the value is destroyed.

However Rust is a at the core a systems programming language so we
need some way to bypass the safety restrictions it provides. This is
where the \textbf{unsafe}~\cite{rustonomicon} keyword comes in. This
keyword provides a way to tell the Rust compiler that a piece of code
upholds the memory safety guarantees even though the former can't
verify  statically that it does. Handling of raw pointers and call to
foreign functions are cases where the use of \textbf{unsafe} is
required. Indeed other languages do not provide any of Rust's
guarantees and do not know about lifetimes. This problem also occurs
when developing kernels~\cite{rust-os}~\cite{redox}.

\section{Concurrency safety}

Rust's type system also provides protection against data-races, a
common problem when using C usually solved through atomic operations
and locks. While it is possible to solve this kind of problems in C,
it requires careful design and is an error-prone task. Rust therefore
allows quicker and safer development of concurrent software. This safety
does not come for free though. In order to guarantee that there is no
data-race anywhere, Rust only allows either one mutable reference or
any number of immutable reference to any given variable. By doing
that, the compiler can ensure that the variable will only be read
concurrently or written exclusively.

Even though we speak of concurrency safety, Rust does not provide
protection against deadlocks. This is far too complex a task for
static analysis in a compiler since it requires costly call graph
analysis and sometimes produces false
positives~\cite{deadlock-detection}.

However, the Rust standard library provides safe ways to mutate shared
variables. As an example the RwLock standard library structure allows
creating mutable reference to a value from an immutable RwLock instance.

\section{Performance}

Being a systems programming language, performance is critical to
Rust. All the aforementioned guarantees are enforced statically so
they do not incur any performance penalty. Removing pointer aliasing
from the language also allows the compiler to do many optimizations
that are not feasible safely in C or C++. As to the performance of the
generated assembly code, Rust utilizes the LLVM compiler
infrastructure and thus benefits from all the work that was done
there. Rust also features what the Rust developers call 0-cost
abstraction. Idiomatic Rust makes heavy use of closures, and that
would incur, in some languages, a slight degradation of
performance. But, in Rust, closure can be, most of the time, inlined
thus making as if the programmer wrote the concrete code while
allowing the interface to remain abstract.

In this regard, Rust combines the best of both worlds, high-level
readable and abstract code while also having performance close to that
of other systems programming languages.

%%% Local Variables:
%%% mode: latex
%%% TeX-master: "master"
%%% End:


\chapter{Design}

\include*{design}

\chapter{Implementation}

\include*{impl}

\chapter{Evaluation}

\include*{eval}

\chapter{Further work}

\include*{work}

\chapter{Conclusion}

We hope to have provided an easy to use and clear API to build upon
our system. We also think that we have proved that Rust is a valid
language to build kernel bypass networking in. \todo{speak about the
  actual performance}

\pagenumbering{gobble}
\bibliographystyle{plain}
\bibliography{master}{}

\end{document}
%%% Local Variables:
%%% mode: latex
%%% TeX-master: t
%%% End:
